\section{Introduction}
\label{sec:intro}

Injection vulnerabilities arise when untrusted data crosses a language boundary
and is misinterpreted as structure rather than as inert content.
A quote can close a SQL string literal, an angle bracket can open an HTML tag,
and a semicolon can start a new shell command.
Despite decades of mitigations (parameterized queries, context-sensitive
output encoding, allowlisting), injection continues to dominate vulnerability
disclosures.
Each defense is specific to one host language and one syntactic context, and
each is applied by hand at the point of use.

The \emph{matchertext} discipline~\cite{matchertext} proposes a single,
language-independent structural rule: the ASCII matcher pairs \verb|()|,
\verb|[]|, and \verb|{}| must always match.
Under this rule, a host language can delimit an untrusted value with a matcher
pair and read it to matcher balance.
To escape its slot, the value would have to emit an unmatched matcher.
Such a matcher is not valid matchertext, and is escaped on the way in.
The canonical injection \verb|' OR '1'='1| illustrates the mechanism: hosted as
\verb|WHERE name = [| \emph{value} \verb|]|, the quote that would have closed
the string literal is now ordinary text inside the matched region, and the
tautology cannot break out.

This paper asks what a matchertext-based defense would guarantee, under what
assumptions, and how much of the real injection landscape it could reach.
We make four contributions:
\begin{itemize}
\item	A security reading of the matchertext discipline
	(\cref{sec:relevance}): why a rule designed for verbatim embedding
	bears on injection, and what would have to hold for it to help.
\item	A threat model (\cref{sec:threat}) stating the guarantee as an
	\emph{inertness} invariant, distinguished from the weaker closure
	theorem and reduced to three named premises; the load-bearing boundary
	premise is discharged by a machine-checked Lean proof
	(\cref{sec:appendix:lean}).
\item	A structural argument (\cref{sec:resist}) that matcher-delimiting
	resists injection across SQL, HTML, LDAP, and PDF, with the options for
	fitting the discipline into existing languages and a comparison with
	deployed defenses.
\item	An empirical measurement of its structural reach (\cref{sec:prevent})
	over a reproducible, snapshot-pinned corpus (\cref{sec:corpus})
	unifying the CVE record with normalized enrichment and proof-of-concept
	databases, classified by weakness family and embedded \emph{syntax}
	(\cref{sec:results}).
\end{itemize}

Together these answer whether one uniform check could replace the collection of
per-host defenses injection is fought with today, and how much of the recorded
landscape such a check would reach.
Both figures below assume a matchertext-aware host and characterize attack
structure rather than a deployed defense. The payload-backed and projected
figures are reported separately, each over its own subset.
